\documentclass[9pt,aspectratio=169,xcolor=table]{beamer}

\usepackage[sfdefault]{roboto}
\usepackage[utf8]{inputenc}
\usepackage[T1]{fontenc}

\usepackage{styles/fluxmacros}
\usefolder{styles}
\usetheme[style=red]{flux}

\definecolor{primaryLight}{HTML}{9B2242}
\definecolor{primary}{HTML}{7B1E3A}
\definecolor{secondary}{HTML}{3A5A6B}
\definecolor{tertiary}{HTML}{8C6A3F}

\usepackage{booktabs}
\usepackage{colortbl}
\usepackage{ragged2e}
\usepackage{amssymb}
\usepackage{multirow}
\usepackage{tikz}
\usetikzlibrary{arrows.meta, positioning, fit, backgrounds, calc, shapes.geometric}
\setbeamertemplate{itemize items}[square]
\usepackage[scaled=.92]{beramono}
\usepackage{listings}
\usepackage{array}
\definecolor{stripe}{gray}{0.94}
\newcommand{\striped}{\rowcolors{1}{white}{stripe}}

\definecolor{codebg}{HTML}{FBF2F4}
\definecolor{codegreen}{rgb}{0,0.45,0}
\definecolor{codestring}{rgb}{0.58,0.1,0.2}
\lstdefinestyle{skpython}{
    basicstyle=\ttfamily\footnotesize,
    keywordstyle=\color{primary}\bfseries,
    commentstyle=\itshape\color{codegreen},
    stringstyle=\ttfamily\color{codestring},
    backgroundcolor=\color{codebg},
    breaklines=true,
    keepspaces=true,
    showspaces=false,
    showstringspaces=false,
    frame=single,
    rulecolor=\color{primary!40},
    numbers=none,
}
\lstdefinestyle{skoutput}{
    basicstyle=\ttfamily\footnotesize,
    backgroundcolor=\color{secondary!8},
    breaklines=true,
    keepspaces=true,
    showspaces=false,
    showstringspaces=false,
    frame=single,
    rulecolor=\color{secondary!50},
    numbers=none,
}
\lstset{language=Python, style=skpython}

\definecolor{cycfill}{HTML}{F3DCE4}
\definecolor{cycline}{HTML}{7B1E3A}
\definecolor{cycdofill}{HTML}{E8EEF0}
\definecolor{cycdoline}{HTML}{3A5A6B}
\tikzset{
  cycstep/.style = {font=\sffamily\small, align=center, rounded corners=2pt,
                     draw=cycline, fill=cycfill, line width=0.7pt,
                     minimum height=7mm, minimum width=20mm, inner sep=3pt},
  cycdo/.style   = {cycstep, fill=cycdofill, draw=cycdoline},
  flow/.style    = {-{Stealth[length=2mm]}, thick, draw=cycline},
}

\AtBeginSection[]
{
{
    \setbeamercolor{background canvas}{bg=primaryLight!6}
    \begin{frame}[plain]
        \frametitle{Table of Contents}
        \tableofcontents[currentsection]
    \end{frame}
}
}

\title{Functions and Reusable Calculations}
\subtitle{\emph{Engineering Computation with Python} --- Chapter 5}
\author{Mun\'{i}m Zabidi}
\institute{\color{primary}Faculty of Electrical Engineering, Universiti Teknologi Malaysia}
\date{\today}
\titlegraphic{assets/logoutmwhite.png}

\begin{document}

\titlepage

\begin{frame}{Table of Contents}
    \tableofcontents
\end{frame}

%===============================================================================
\section{What Is a Function?}
%===============================================================================

\begin{frame}{Why Functions?}
    \leavevmode
    So far every program has been written as one continuous script. That
    works while a calculation is short, but engineering calculations are
    \alert{repeated} --- the same conversion applied to twenty readings,
    the same formula checked at a dozen operating points.

    \vspace{3mm}
    \begin{block}{A function}
        Packages a calculation once, gives it a name, and lets the rest
        of the program call it as often as needed.
    \end{block}

    \vspace{2mm}
    \begin{itemize}\small
        \item Makes code cleaner and easier to read
        \item Improves reusability and maintainability
        \item Hides complex implementation details
    \end{itemize}
\end{frame}

\begin{frame}{Function vs. Script}
    \leavevmode
    \small
    \striped
    \begin{center}
    \begin{tabular}{lll}
        \toprule
        \textbf{Aspect} & \textbf{Script} & \textbf{Function} \\
        \midrule
        Usage & Executes commands directly & Defines a reusable task \\
        Input/output & Global variables, user input & Arguments, return values \\
        Scope & Uses global variables & Uses local variables \\
        \bottomrule
    \end{tabular}
    \end{center}

    \vspace{4mm}
    \begin{alertblock}{The key shift}
        A script runs top to bottom, once. A function is written once
        and \alert{called} as many times as needed, with different
        inputs each time.
    \end{alertblock}
\end{frame}

\begin{frame}{Types of Functions}
    \leavevmode
    \small
    \striped
    \begin{center}
    \begin{tabular}{p{2.6cm}p{4.3cm}p{3.3cm}}
        \toprule
        \textbf{Type} & \textbf{Description} & \textbf{Example} \\
        \midrule
        Built-in & Predefined in Python & \texttt{print()}, \texttt{len()} \\
        Library & From external modules & \texttt{math.sqrt()} \\
        User-defined & Created by the programmer & \texttt{def greet():} \\
        Nested (local) & Defined inside another function & see below \\
        \bottomrule
    \end{tabular}
    \end{center}

    \vspace{3mm}
    {\small This chapter covers all four, with the focus on
    \alert{writing your own}.}
\end{frame}

%===============================================================================
\section{Pre-Defined Functions}
%===============================================================================

\begin{frame}[fragile]{Built-in Functions: Always Available}
    \leavevmode
    Built-in functions need \alert{no import} at all:

\begin{lstlisting}
readings = [12.0, 11.5, 13.2, 10.8]   # volts

print(len(readings))                   # 4
print(max(readings))                   # 13.2
print(min(readings))                   # 10.8
print(sum(readings) / len(readings))   # 11.875 -- the mean
\end{lstlisting}

    \vspace{1mm}
    {\small \texttt{sum(x)/len(x)} is worth remembering: an average is
    nothing more than a total divided by a count. Chapter 14 replaces
    this with \texttt{np.mean()}.}
\end{frame}

\begin{frame}[fragile]{A Surprise in round()}
    \leavevmode
    Python rounds a value ending in exactly $.5$ to the nearest
    \alert{even} number, not always upward:

\begin{lstlisting}
print(round(2.5))    # 2, not 3
print(round(3.5))    # 4
\end{lstlisting}

    \vspace{3mm}
    \begin{block}{Banker's rounding}
        This exists so that rounding a large set of measurements does
        not accumulate a systematic upward bias. Correct behaviour, not
        a defect --- but it surprises almost everyone the first time.
    \end{block}
\end{frame}

\begin{frame}[fragile]{Library Functions}
    \leavevmode
    Library functions become available through \alert{importing}:

\begin{lstlisting}
import math
print(math.sqrt(16))           # 4.0

import numpy as np
array = np.array([1, 2, 3])

import random
choice = random.choice(['apple', 'banana', 'cherry'])
\end{lstlisting}

    \vspace{1mm}
    {\small Every library used in this course --- \texttt{math},
    \texttt{numpy}, \texttt{pandas}, \texttt{random} --- follows this
    same \texttt{import} pattern.}
\end{frame}

%===============================================================================
\section{Writing Your Own Functions}
%===============================================================================

\begin{frame}[fragile]{Function Syntax}
    \leavevmode
    A user-defined function uses \texttt{def}, a name, parameters, a
    colon, and an indented body. \texttt{return} sends a result back.

\begin{lstlisting}
def function_name(parameter1, parameter2):
    """Description of what the function does"""
    # Code block
    return result
\end{lstlisting}

\begin{lstlisting}
def add(a, b):
    return a + b

print(add(3, 7))  # 10
\end{lstlisting}
\end{frame}

\begin{frame}{Five Steps to Define a Function}
    \leavevmode
    \begin{enumerate}
        \item \textbf{Define}: \texttt{def} keyword, name, parentheses
        \item \textbf{Parameters} (optional): inputs inside the parentheses
        \item \textbf{Colon \texttt{:}}: starts the function body
        \item \textbf{Body}: indented code
        \item \textbf{Return} (optional): sends back a result
    \end{enumerate}

    \vspace{3mm}
    \begin{block}{Remember}
        A function that never uses \texttt{return} still runs fine ---
        it simply returns \texttt{None} to the caller.
    \end{block}
\end{frame}

\begin{frame}[fragile]{Single Input, Single Output}
    \leavevmode
    The simplest shape: one value in, one value out.

\begin{lstlisting}
def square(number):
    return number ** 2

print(square(4))  # 16
\end{lstlisting}

    \vspace{3mm}
    {\small The four input/output shapes --- single/single,
    single/multiple, multiple/single, multiple/multiple --- are all
    equally valid in Python. The next few slides show each.}
\end{frame}

\begin{frame}[fragile]{Single Input, Multiple Output}
    \leavevmode
    Multiple return values are packed together as a \alert{tuple}:

\begin{lstlisting}
def analyze_number(num):
    is_even = num % 2 == 0
    square = num ** 2
    return is_even, square

even, sq = analyze_number(5)
print(even, sq)
# False 25
\end{lstlisting}

    \vspace{1mm}
    {\small Returning several values at once is idiomatic in Python ---
    and convenient for engineering results.}
\end{frame}

\begin{frame}[fragile]{Multiple Input, Multiple Output}
    \leavevmode
\begin{lstlisting}
def calculate(a, b):
    addition = a + b
    subtraction = a - b
    multiplication = a * b
    division = a / b if b != 0 else None
    return addition, subtraction, multiplication, division

s, d, p, q = calculate(8, 2)
print(s, d, p, q)
# 10 6 16 4.0
\end{lstlisting}

    \vspace{1mm}
    {\small Any combination of inputs and outputs is possible --- choose
    the shape that matches the calculation.}
\end{frame}

\begin{frame}[fragile]{Selective Output Handling}
    \leavevmode
    The caller can choose to use only the outputs it needs:

\begin{lstlisting}
def motion(t):
    distance = [0.0001 * i**2 for i in t]
    acceleration = [0.0025 * i**2 for i in t]
    return distance, acceleration

distance, acceleration = motion(t)   # both
distance = motion(t)[0]              # index into the tuple
distance, _ = motion(t)              # discard the 2nd value
_, acceleration = motion(t)          # discard the 1st value
\end{lstlisting}

    \vspace{1mm}
    {\small The underscore \texttt{\_} is a common convention for
    ``I don't need this value.''}
\end{frame}

\begin{frame}[fragile]{Local (Nested) Functions}
    \leavevmode
    A function defined \alert{inside} another function is a local
    function --- accessible only from within the enclosing function.

\begin{lstlisting}
def process_order(quantity, price_per_item):
    def calculate_total(qty, price):      # local function
        total = qty * price
        if total > 100:
            return total * 0.9  # 10% discount
        return total

    total_price = calculate_total(quantity, price_per_item)
    return f"Total Price: ${total_price:.2f}"

print(process_order(10, 15))
# Total Price: $135.00
\end{lstlisting}
\end{frame}

\begin{frame}[fragile]{Multiple Local Functions and Default Arguments}
    \leavevmode
\begin{lstlisting}
def process_order(quantity, price_per_item):
    def calculate_total(qty, price):
        total = qty * price
        return total * 0.9 if total > 100 else total

    def apply_tax(amount, tax_rate=0.05):  # default argument
        return amount * (1 + tax_rate)

    total_price = apply_tax(calculate_total(quantity, price_per_item))
    return f"Total after tax: {total_price}"
\end{lstlisting}

    \vspace{1mm}
    {\small \texttt{apply\_tax} has a \alert{default argument}: if
    \texttt{tax\_rate} is not supplied, it defaults to \texttt{0.05}.}
\end{frame}

\begin{frame}[fragile]{Documenting a Function: Docstrings}
    \leavevmode
    A \alert{docstring} documents what a function does, what it expects,
    and what it returns.

\begin{lstlisting}
def calculate_circle_area(radius):
    """
    Calculate the area of a circle given its radius.

    Parameters:
        radius (float): The radius of the circle.
    Returns:
        float: The area of the circle.
    """
    return math.pi * radius ** 2

help(calculate_circle_area)
\end{lstlisting}

    \vspace{1mm}
    {\small The difference between a function a colleague can use, and
    one they must reverse-engineer.}
\end{frame}

%===============================================================================
\section{Variable Scope}
%===============================================================================

\begin{frame}{What Is Scope?}
    \leavevmode
    Every variable lives within a certain \alert{scope} --- the region
    of code where it can be seen or accessed.

    \vspace{3mm}
    \begin{columns}[t]
        \begin{column}{.5\textwidth}
            \begin{block}{Local}
                \small Defined inside a function. Only works there.
                Deleted when the function finishes.
            \end{block}
        \end{column}
        \begin{column}{.5\textwidth}
            \begin{block}{Global}
                \small Defined outside all functions. Accessible
                throughout the program, including inside functions.
            \end{block}
        \end{column}
    \end{columns}

    \vspace{3mm}
    {\small This distinction causes some of the most confusing bugs
    beginners encounter --- worth understanding carefully.}
\end{frame}

\begin{frame}[fragile]{Local Variables Disappear}
    \leavevmode
\begin{lstlisting}
def greet():
    name = "Alice"  # local variable
    print(name)

greet()
print(name)  # NameError: name is not defined
\end{lstlisting}

    \vspace{2mm}
    \begin{alertblock}{What happened}
        \texttt{name} only exists while \texttt{greet()} is running.
        Once the function finishes, it is \alert{deleted from memory}
        --- the outer \texttt{print(name)} has nothing to find.
    \end{alertblock}
\end{frame}

\begin{frame}[fragile]{Global Variables Are Visible Everywhere}
    \leavevmode
\begin{lstlisting}
name = "Bob"  # global variable

def greet():
    print(name)   # Python finds it in the global scope

greet()        # Bob
print(name)    # Bob
\end{lstlisting}

    \vspace{2mm}
    \begin{block}{What happened}
        There is no local \texttt{name} inside \texttt{greet()}, so
        Python automatically looks outward and finds the global one.
    \end{block}
\end{frame}

\begin{frame}[fragile]{The Same Name, Two Different Variables}
    \leavevmode
    When a function \alert{assigns} to a name, Python treats that name
    as local for the \emph{entire} function body:

\begin{lstlisting}
count = 0  # global

def increase():
    count = 5  # creates a NEW local variable
    print("Inside:", count)

increase()
print("Outside:", count)
\end{lstlisting}
\begin{lstlisting}[style=skoutput]
Inside: 5
Outside: 0
\end{lstlisting}

    \vspace{1mm}
    {\small The local \texttt{count} \alert{shadows} the global one ---
    but only inside the function.}
\end{frame}

\begin{frame}[fragile]{Modifying a Global: The Wrong Way}
    \leavevmode
\begin{lstlisting}
count = 0  # global variable

def increase():
    count = count + 1  # tries to change count
    print(count)

increase()
\end{lstlisting}
\begin{lstlisting}[style=skoutput]
UnboundLocalError: local variable 'count'
referenced before assignment
\end{lstlisting}

    \vspace{1mm}
    \begin{alertblock}{Why}
        Python sees \texttt{count = ...} and assumes \texttt{count} is
        local for the \alert{whole} function --- then tries to read it
        before it has been created.
    \end{alertblock}
\end{frame}

\begin{frame}[fragile]{Modifying a Global: The Right Way}
    \leavevmode
    Use the \texttt{global} keyword to tell Python which \texttt{count}
    you mean:

\begin{lstlisting}
count = 0  # global variable

def increase():
    global count   # use the global one, not a new local
    count = count + 1
    print(count)

increase()
print(count)
\end{lstlisting}
\begin{lstlisting}[style=skoutput]
1
1
\end{lstlisting}

    \vspace{1mm}
    \begin{alertblock}{Rule of thumb}
        Declare \texttt{global} where you intend to \alert{change} the
        value, not merely read it. Prefer parameters and return values
        where you can.
    \end{alertblock}
\end{frame}

%===============================================================================
\section{Testing Functions}
%===============================================================================

\begin{frame}[fragile]{Testing with assert}
    \leavevmode
    A test is simply a call whose result you already know, checked with
    \texttt{assert}:

\begin{lstlisting}
def add(a, b):
    return a + b

assert add(2, 3) == 5
assert add(-1, 1) == 0
assert add(0, 0) == 0
\end{lstlisting}

    \vspace{2mm}
    \begin{block}{Pass vs. fail}
        If all assertions hold, nothing is printed. If one fails,
        Python stops and tells you \alert{exactly which case} broke.
    \end{block}
\end{frame}

\begin{frame}{Choosing Good Test Vectors}
    \leavevmode
    \small
    \striped
    \begin{center}
    \begin{tabular}{p{2.6cm}p{3.0cm}p{3.6cm}}
        \toprule
        \textbf{Type} & \textbf{Example} & \textbf{Why important} \\
        \midrule
        Typical case & \texttt{add(2, 3)} & Normal usage \\
        Edge case & \texttt{add(0, 0)} & Limits or boundaries \\
        Negative case & \texttt{add(-5, 10)} & Behaviour with negatives \\
        Unusual input & \texttt{add(999999, 1)} & Large values \\
        \bottomrule
    \end{tabular}
    \end{center}

    \vspace{3mm}
    {\small Good test vectors cover a \alert{range} of situations, not
    just the case you expect to work.}
\end{frame}

\begin{frame}[fragile]{Worked Example: An Off-by-One Bug}
    \leavevmode
    Children aged 12 and under get a discount. This implementation has
    a subtle bug:

\begin{lstlisting}
def legoland_price(age, base_price=50.0):
    if age < 12:   # BUG: 12-year-olds should ALSO get the discount
        return base_price * 0.5
    return base_price

assert legoland_price(11) == 25.0           # passes
assert legoland_price(12) == 25.0           # FAILS!
assert legoland_price(13) == 50.0           # passes
\end{lstlisting}

    \vspace{1mm}
    \begin{alertblock}{The fix}
        Change \texttt{age < 12} to \texttt{age <= 12}. The failing
        assertion at the \alert{boundary} (age 12) is exactly what
        revealed the bug.
    \end{alertblock}
\end{frame}

%===============================================================================
\section{Summary}
%===============================================================================

\begin{frame}{What You Should Take Away}
    \leavevmode
    \small
    \begin{itemize}
        \item A \textbf{function} packages a calculation under a name so
              it can be called repeatedly; a \textbf{script} executes
              once.
        \item Python has \textbf{built-in} functions (\texttt{print()},
              \texttt{len()}), \textbf{library} functions
              (\texttt{math.sqrt()}), and \textbf{user-defined}
              functions you write with \texttt{def}.
        \item A function can take any number of inputs and return any
              number of outputs --- returning several values as a tuple
              is idiomatic Python.
        \item Variables inside a function are \textbf{local} and vanish
              when it returns; this isolation prevents a function from
              silently corrupting the rest of the program.
    \end{itemize}
\end{frame}

\begin{frame}{What You Should Take Away (continued)}
    \leavevmode
    \small
    \begin{itemize}
        \item \textbf{Global} variables are visible everywhere, but
              assigning to a name inside a function creates a
              \emph{new local} variable unless you declare
              \texttt{global} first.
        \item A \textbf{docstring} documents a function for the next
              person (often yourself) who has to use it.
        \item Test with \texttt{assert} and good \textbf{test vectors}
              --- typical, edge, negative, and unusual cases, especially
              at boundaries.
    \end{itemize}

    \vspace{3mm}
    \begin{alertblock}{Looking ahead}
        Every function here took one value in and returned one value
        out. A real instrument logs hundreds of readings --- Chapter 6
        introduces the \alert{array}, which computes on many values at
        once.
    \end{alertblock}

    \vspace{3mm}
    \begin{center}
        \large\color{primary}\textbf{Next: Arrays and Vectorized Computation}
    \end{center}
\end{frame}

\end{document}
